% Section 8.1, from lectures/L08/appendix/a_arithmetic_and_flags.md.
%
% Typesetting only: in the subtraction example of 8.1.4 the comments start four columns further
% left than in the lecture, so that the longest line fits the listing's width.

\needspace{10\baselineskip}
\section{Aritmetik och flaggor}
\label{c8:app:a}

\subsection{Addition i ett register}
\label{c8:a1}
Instruktionen \code{add} adderar två register och lägger summan i det första:

\begin{asmcode}
    ldi r16, 25
    ldi r17, 17
    add r16, r17                    ; r16 = r16 + r17 = 42. r17 is unchanged.
\end{asmcode}

\noindent Processorn räknar binärt, precis som i \secref{c1:app:a}: bit för bit från höger, med en
minnessiffra, \textbf{carry}, när en kolumn blir 2 eller mer.

\begin{textdrawing}
   carry     0 0 1 0 0 0 1 0  (in to each column)
   25    =   0 0 0 1 1 0 0 1
 + 17    =   0 0 0 1 0 0 0 1
   ---------------------------
   42    =   0 0 1 0 1 0 1 0  =  0x2A
\end{textdrawing}

\noindent Bit~0: 1 + 1 = 10 binärt, alltså 0 och en etta i minne till bit~1. Bit~4: 1 + 1 = 10
igen, en etta i minne till bit~5. Resten av kolumnerna blir 0 eller 1 utan minne.

Två instruktioner till adderar:
\begin{itemize}
  \item \textbf{\code{inc r16}} adderar 1. Den har redan dykt upp i kapitel~\ref{c7:ch}.
  \item \textbf{\code{adc r16, r17}} adderar \code{r17} \textbf{och carryflaggan} till \code{r16}.
    Den behövs för tal större än en byte, och tas upp i kapitel~\ref{c10:ch}.
\end{itemize}

Det finns \textbf{ingen instruktion som adderar en konstant}, ingen \code{addi}. Hur man ändå gör
det står i \secref{c8:a4}.

\subsection{När summan inte får plats}
\label{c8:a2}
En byte har åtta bitar och rymmer talen 0 till 255. Vad händer med 200 + 100 = 300?

\begin{textdrawing}
   200   =     1 1 0 0 1 0 0 0
 + 100   =     0 1 1 0 0 1 0 0
   -----------------------------
   300   =   1 0 0 1 0 1 1 0 0
             ^
             nionde biten: hamnar i carryflaggan C
\end{textdrawing}

\noindent Summan behöver nio bitar. Registret behåller de åtta lägsta, \code{0010 1100} =
\textbf{44}, och den nionde biten, minnessiffran ut ur bit~7, hamnar i \textbf{carryflaggan
\code{C}} i statusregistret. Så \code{C = 1} efter additionen betyder: \enquote{resultatet fick inte
plats, det riktiga svaret är 256 större}. 44 + 256 = 300.

Processorn stannar inte, och säger inte ifrån. Det är programmets sak att titta på \code{C} om det
spelar roll, och det är just det flaggorna är till för (\secref{c8:a6}).

\subsection{Aritmetisk rundgång}
\label{c8:a3}
En byte fungerar som mätaren i en bil: efter 999\,999 kommer 000\,000. Efter 255 kommer 0, och före
0 kommer 255. Det kallas \textbf{aritmetisk rundgång}.

\bookfigure{l08_number_circle}{De 256 värdena i en byte som en urtavla, utan tecken utanför och
med tecken innanför.}{c8:fig:number_circle}

Tänk dig talen på en urtavla. \code{inc} flyttar ett steg medurs, \code{dec} ett steg moturs. Från
255 tar \code{inc} dig till 0, och från 0 tar \code{dec} dig till 255.

\textbf{Här finns en fälla.} \code{inc} och \code{dec} ändrar \textbf{inte} carryflaggan. När
\code{inc} slår om från 255 till 0 sätts \code{Z}, eftersom resultatet blev noll, men \code{C}
lämnas som den var. Vill man veta om en ökning slog om måste man antingen använda \code{add} med ett
register som innehåller 1, som sätter \code{C}, eller titta på \code{Z} efter \code{inc}.

\begin{narrowtable}{@{}lcccc@{}}
  \thead{Instruktion} & \thead{Från} & \thead{Till} & \thead{Z} & \thead{C} \\
  \midrule
  \code{inc r16} & 255 & 0 & 1 & oförändrad \\
  \code{add r16, r17} med \code{r17} = 1 & 255 & 0 & 1 & 1 \\
  \code{dec r16} & 0 & 255 & 0 & oförändrad \\
\end{narrowtable}

\noindent Varför är \code{inc} och \code{dec} gjorda så? Därför att de oftast används som räknare i
loopar, och då vill man kunna räkna utan att förstöra \code{C}, som kanske håller något viktigt.
Det blir tydligt i kapitel~\ref{c10:ch}, där \code{C} bär minnessiffran mellan två byte.

\subsection{Subtraktion}
\label{c8:a4}
Subtraktion finns i tre varianter:

\begin{booktable}{@{}p{6em}p{7.5em}L@{}}
  \thead{Instruktion} & \thead{Exempel} & \thead{Gör} \\
  \midrule
  \code{sub} & \code{sub r16, r17} & \code{r16 = r16 - r17} \\
  \code{subi} & \code{subi r16, 5} & \code{r16 = r16 - 5}, med en konstant. Bara
    \code{r16}--\code{r31}. \\
  \code{dec} & \code{dec r16} & \code{r16 = r16 - 1} \\
\end{booktable}

\noindent Om det man drar ifrån är större än det man drar ifrån det från, får man \textbf{låna} från
en tänkt nionde bit. Då sätts \textbf{\code{C = 1}}, som här betyder \emph{lån} (\emph{borrow}) i
stället för minnessiffra:

\begin{asmcode}
    ldi r16, 20
    ldi r17, 50
    sub r16, r17                ; 20 - 50: r16 = 226 = 0xE2, and C = 1 (a borrow)
\end{asmcode}

\noindent 226 är $256 + 20 - 50$. Rundgången gäller åt båda hållen.

\subsubsection{Att addera en konstant med \code{subi}}
Eftersom det inte finns någon \code{addi} adderar man en konstant genom att \textbf{subtrahera dess
negativa värde}:

\begin{asmcode}
    subi r16, -5                    ; r16 = r16 - (-5) = r16 + 5
\end{asmcode}

\noindent Det fungerar på grund av rundgången: $-5$ lagras som 251 (\secref{c8:a5}), och att dra
ifrån 251 är detsamma som att lägga till 5 och slå om ett varv. Tänk på att flaggorna då blir en
subtraktions flaggor: \code{C} betyder lån, inte minnessiffra, så \code{C} blir 1 när summan
\textbf{inte} slår om, tvärtom mot \code{add}. Använd \code{subi} med negativ konstant för att
räkna, inte för att testa carry.

\subsection{Tvåkomplement: negativa tal}
\label{c8:a5}
En byte är bara åtta bitar. Om de betyder ett tal utan tecken, 0 till 255, eller ett tal med
tecken, $-128$ till 127, bestämmer \textbf{programmet}, inte processorn. Samma bitar kan läsas på
två sätt:

\begin{narrowtable}{@{}cccc@{}}
  \thead{Bitar} & \thead{Hex} & \thead{Utan tecken} & \thead{Med tecken} \\
  \midrule
  \code{0000 0000} & 0x00 & 0 & 0 \\
  \code{0000 0001} & 0x01 & 1 & 1 \\
  \code{0111 1111} & 0x7F & 127 & 127 \\
  \code{1000 0000} & 0x80 & 128 & $-128$ \\
  \code{1111 1011} & 0xFB & 251 & $-5$ \\
  \code{1111 1111} & 0xFF & 255 & $-1$ \\
\end{narrowtable}

\noindent Sättet att skriva negativa tal kallas \textbf{tvåkomplement}. Det har tre egenskaper som
gör det till det alla processorer använder:
\begin{itemize}
  \item \textbf{Bit~7 är teckenbiten.} Är den 1 är talet negativt, när byten läses med tecken.
  \item \textbf{$-x$ är $256 - x$.} $-1$ är 255, $-5$ är 251 och $-128$ är 128. Det är rundgången
    igen: ett steg moturs från 0 är både 255 och $-1$.
  \item \textbf{Addition och subtraktion fungerar likadant för båda tolkningarna.} \code{add}
    behöver inte veta om talen har tecken: 0xFB + 0x07 = 0x02 är både $251 + 7 = 258 - 256 = 2$
    och $-5 + 7 = 2$.
\end{itemize}

\subsubsection{Att bilda ett negativt tal}
Så räknar du ut tvåkomplementet, det negativa talet, för hand: \textbf{invertera alla bitar och
lägg till 1.}

\begin{textdrawing}
   5            =  0000 0101
   invertera    =  1111 1010
   lägg till 1  =  1111 1011  =  0xFB  =  -5
\end{textdrawing}

\noindent Instruktionen \textbf{\code{neg r16}} gör exakt detta: \code{r16 = 0 - r16}.

Det är samma sak Windows miniräknare visar i läget Programmerare när den visar ett negativt tal som
\code{FFFF FFFF FFFF FFFB}: det är $-5$ i tvåkomplement med 64 bitar i stället för 8.

\subsection{Flaggorna}
\label{c8:a6}
Efter varje aritmetisk eller logisk instruktion skriver processorn in några fakta om resultatet i
statusregistret \code{SREG}, en bit per faktum. Bitarna kallas \textbf{flaggor}.

\bookfigure{avr_sreg}{Statusregistret SREG och dess flaggor.}{c8:fig:sreg}

\begin{booktable}{@{}p{3.6em}p{5.6em}L@{}}
  \thead{Flagga} & \thead{Namn} & \thead{Sätts till 1 när} \\
  \midrule
  \code{C} & carry & det blev en minnessiffra ut ur bit~7 vid addition, eller ett lån vid
    subtraktion \\
  \code{Z} & zero & resultatet blev noll \\
  \code{N} & negative & bit~7 i resultatet är 1, alltså negativt om talet har tecken \\
  \code{V} & overflow & resultatet blev fel \textbf{med tecken}: två positiva tal gav ett
    negativt, eller tvärtom \\
  \code{S} & sign & \code{N} xor \code{V}: det rätta tecknet, även när \code{V} säger att
    resultatet slog om \\
  \code{H} & half carry & det blev en minnessiffra ut ur bit~3; används sällan \\
\end{booktable}

\noindent De två sista bitarna, \code{I} och \code{T}, är inte resultatflaggor och används inte i
boken.

\textbf{\code{C} och \code{V} svarar på samma fråga för de två tolkningarna.} \code{C = 1} betyder
att resultatet inte fick plats om talen är \textbf{utan} tecken. \code{V = 1} betyder att
resultatet inte fick plats om talen är \textbf{med} tecken. Processorn vet inte vilken tolkning
programmet använder, så den räknar ut båda, och programmet tittar på den flagga som passar.

\textbf{Inte alla instruktioner ändrar alla flaggor.} \code{ldi}, \code{mov}, \code{out} och
hoppen ändrar inga flaggor alls. \code{inc} och \code{dec} ändrar \code{Z}, \code{N}, \code{V} och
\code{S}, men inte \code{C}. Vilka flaggor varje instruktion ändrar står i kolumnen \emph{Flaggor}
på referensbladet i bilaga~\ref{app:avr}. Det är värt att titta där när ett program beter sig
konstigt: en flagga som man trodde var satt kan ha skrivits över av instruktionen före.

\subsection{Fyra genomarbetade exempel}
\label{c8:a7}
Programmet \code{flags_demo.asm} gör de fyra uträkningarna nedan. Stega det och jämför SREG i
Processor Status efter varje uträkning.

\asmfile{lectures/L08/examples/flags_demo.asm}

\subsubsection{0x7F + 0x01 = 0x80}

\begin{textdrawing}
   0111 1111    127  eller  +127
 + 0000 0001      1  eller    +1
   ---------
   1000 0000    128  eller  -128
\end{textdrawing}

\begin{itemize}
  \item \code{C = 0}: ingen minnessiffra ut ur bit~7. Utan tecken är 128 rätt.
  \item \code{Z = 0}: resultatet är inte noll.
  \item \code{N = 1}: bit~7 är 1.
  \item \code{V = 1}: två positiva tal gav ett negativt. Med tecken är svaret fel, $+127 + 1$ får
    inte plats.
  \item \code{S = N xor V = 0}: det rätta svaret, 128, är positivt.
  \item \code{H = 1}: bit~3 + bit~3 gav en minnessiffra, 1111 + 0001 i den låga halvan.
\end{itemize}

\subsubsection{0xFF + 0x01 = 0x00}
\begin{itemize}
  \item \code{C = 1}: 255 + 1 = 256 får inte plats utan tecken.
  \item \code{Z = 1}: resultatet är noll.
  \item \code{N = 0}, \code{V = 0}: med tecken är det $-1 + 1 = 0$, vilket är rätt.
  \item \code{S = 0}, \code{H = 1}.
\end{itemize}

\subsubsection{0x80 + 0x80 = 0x00}
\begin{itemize}
  \item \code{C = 1}: 128 + 128 = 256 får inte plats utan tecken.
  \item \code{Z = 1}: resultatet är noll.
  \item \code{V = 1}: med tecken är det $-128 + (-128) = -256$, som inte heller får plats. Två
    negativa tal gav noll, och det är fel.
  \item \code{N = 0}, \code{S = N xor V = 1}: det rätta svaret, $-256$, är negativt. \code{H = 0}.
\end{itemize}

\subsubsection{0x05 -- 0x07 = 0xFE}
\begin{itemize}
  \item \code{C = 1}: $5 - 7$ kräver ett lån. Utan tecken blir det 254.
  \item \code{Z = 0}, \code{N = 1}: bit~7 är 1.
  \item \code{V = 0}: med tecken är $5 - 7 = -2$, och 0xFE \textbf{är} $-2$. Svaret är rätt.
  \item \code{S = 1}: resultatet är negativt. \code{H = 1}: den låga halvan, $5 - 7$, behövde
    också låna.
\end{itemize}

Lägg märke till det sista exemplet: \code{C = 1} men \code{V = 0}. Utan tecken är svaret
\enquote{fel} (det skulle ha varit $-2$, som inte finns bland talen 0--255), men med tecken är det
rätt. Samma bitar, två tolkningar, två flaggor.
